\section{Conclusion and Future Work}\label{sec:conclusion}
This work presents SFRRT*, a sampling-based motion planner for the spill-free transport of open-top, liquid-filled containers, particularly in cluttered scenes. SFRRT* distinguishes itself by utilizing implicitly learned spillage dynamics via a transformer-based machine-learning approach. This outperforms existing approaches by increasing the success rate in cluttered scenes. Specifically, SFRRT* enables motions across all axes of end effector orientations using the learned model instead of simplified liquid models. It also permits the container to tilt to its maximum tilt angle, thus enabling flexible maneuvering to avoid obstacles. It also does not need visual monitoring of the liquid, which could get obstructed in cluttered scenes. 
% This is because it relies solely on the container motion to ensure spill-free trajectory generation.


However, like many machine learning models, our approach faces a common limitation— the need for a substantial amount of training data. Moreover, containers are approximated as frustums of cones to estimate the maximum tilt angles. This could lead to conservative estimates for containers with more complex curvy shapes. An alternative would be training a model on image inputs to predict tilt angles more precisely. Lastly, circular upside-down motions (loop-de-loop) can avoid spilling if executed with adequate speed. Still, SFRRT* cannot support generating them because it constrains container orientations to be less than the maximum tilt angle. This limitation stems from the fact that SFC is trained on data collected from humans, who typically exhibit caution when handling liquids. Consequently, the model lacks training data to classify loop-the-loop motions as spill-free.


For future work, we are interested in improving the Spill-Free Classifier model to estimate spillage probability rather than solely classifying it, thus allowing for a more precise time-parameterized trajectory. Furthermore, the path planning phase currently only constrains the orientation range within the maximum tilt angle. However, if the sampled orientation in the resulting path approaches the maximum tilt angle, the generated trajectory will be very slow. To address this, we plan to analyze how the sampled orientation affects velocity, acceleration, and jerk during path planning. We aim to generate a time-optimal parameterized path by considering these kinematic factors in the path-planning phase.
Moreover, our approach is versatile and can be expanded to a variety of applications. We plan to extend our method to include spillage dynamics for granular materials such as sand and sugar, to transport a tablespoon of these materials without spillage.
Furthermore, our method is adaptable to different path planners and not solely dependent on RRT*. Moving forward, we plan to evaluate our approach using various planning algorithms and compare success rates and planning times.
Finally, to improve our SFC accuracy and SFRRT* success rate further, we aim to optimize our data collection process by exploring automation through Computational Fluid Dynamics (CFD) frameworks. 
